\documentclass[11pt]{article}
\usepackage[a4paper,margin=1.1in]{geometry}
\usepackage[T1]{fontenc}
\usepackage{lmodern}
\usepackage{amsmath,amssymb,amsthm}
\usepackage{booktabs}
\usepackage{longtable}
\usepackage{array}
\usepackage{microtype}
\usepackage[hidelinks]{hyperref}
\usepackage{natbib}
\usepackage{setspace}
\usepackage{enumitem}
\usepackage{xcolor}

\newtheorem{theorem}{Theorem}[section]
\newtheorem{proposition}[theorem]{Proposition}
\newtheorem{lemma}[theorem]{Lemma}
\newtheorem{corollary}[theorem]{Corollary}
\theoremstyle{definition}
\newtheorem{definition}[theorem]{Definition}
\newtheorem{assumption}[theorem]{Assumption}
\theoremstyle{remark}
\newtheorem{remark}[theorem]{Remark}

% A formal name, as it appears in the Lean development.
\newcommand{\leanunderscore}{\char`\_\hskip 0pt plus 0.4pt\relax}
\newcommand{\lean}[1]{\begingroup\ttfamily\small\let\_\leanunderscore #1\endgroup}
% The Lean name attached to a stated result.
\newcommand{\formal}[1]{\par\noindent\begingroup\footnotesize\raggedright\emph{Formal statement:} \lean{#1}.\par\endgroup}
\newcommand{\Thorn}{\text{\th}}
\newcommand{\R}{\mathbb{R}}
\newcommand{\E}{\mathbb{E}}
\DeclareMathOperator{\supp}{supp}

\onehalfspacing
\tolerance=600 \emergencystretch=1.5em

\title{Uniqueness of the Stationary Equilibrium in Aiyagari (1994):\\
A Machine-Checked Proof for the Log Calibration,\\ and the Frontier for Higher Risk Aversion}
\author{Robert Kirkby\thanks{Victoria University of Wellington. The Lean development described here
is at \url{https://github.com/robertdkirkby/EconomicsInLean}. It was written with the assistance of
an AI coding model (Claude, Anthropic); every result reported was checked by the Lean kernel.
Comments welcome.}}
\date{September 2026\\[0.5em]\small Draft}

\begin{document}
\maketitle

\begin{abstract}
\noindent
\citet{Aiyagari1994} computed the stationary equilibrium of his heterogeneous-agent economy
numerically. We prove that at his logarithmic calibration
($\mu=1$) it is: for any of his labour-endowment processes, and for any finite Markov earnings
process with monotone transitions, at most one interest rate in $(-\delta,\lambda)$ clears the
capital market, where $\lambda=1/\beta-1$ is the rate of time preference; no equilibrium exists
above $\lambda$ or below $-4\%$; and the bisection he used is certified to converge to the rate.
Every step is checked by the Lean 4 proof assistant against the Mathlib library. The economics is Light's (2020) chain, saving rises with the rate, so aggregate capital
does, so supply meets a decreasing demand once, but at $\beta=0.96$ the published constants fail; what
succeeds is the Euler inequality at the borrowing constraint, a
minimal marginal propensity to consume, a Doeblin atom, and homogeneity of the stationary
distribution in the wage. For $\mu\in\{3,5\}$ we locate the obstruction exactly: a condition on
the response of consumption to the rate that fails for the rich. A tighter argument weakens it to
one that holds numerically at every wealth level with an explicit margin, but that remains
unproved. For \citet{Huggett1993} the obstruction is the debtors, and we bound their effect.
\end{abstract}

\noindent\emph{JEL:} C63, D52, E21. \quad \emph{Keywords:} Bewley models, uniqueness, stationary
equilibrium, formal verification, Lean.

\section{Introduction}

The stationary equilibrium of \citet{Aiyagari1994} is a number: the interest rate at which the
capital that a continuum of households accumulates against uninsurable earnings risk equals the
capital a competitive firm demands. Aiyagari computed it by bisection on the rate and reported it,
for three degrees of relative risk aversion $\mu\in\{1,3,5\}$ and eight earnings processes, in his
Table II. He did not ask whether the number was unique, and for three decades the question was
treated as settled by the picture: capital supply is drawn rising in the rate, demand falling, and
the curves cross once.

The picture is not a theorem. \citet{Acikgoz2018} proved existence and gave conditions for
uniqueness that do not cover Aiyagari's parameters; \citet{Light2020} proved uniqueness for
relative risk aversion at most one, again under conditions on the discount factor and the earnings
process that his own statement does not verify at $\beta=0.96$; \citet{Light2022} put the argument
in a general mean-field framework; \citet{Kirkby2019} computed the Aiyagari equilibrium with a discretisation proved to converge to
the true solution, found numerically that it is almost certainly unique, and recorded that
uniqueness remained open theoretically; \citet{YoungWalsh2025} showed by computation that at high
risk aversion the supply curve can bend backward and the equilibrium can fail to be unique. The survey
of \citet{TodaWalsh2024} lists uniqueness in the Aiyagari model with its actual calibration as
open.

This paper closes it for $\mu=1$ and locates the obstruction for $\mu\in\{3,5\}$, and it does so in
a form that is new to economics: every result is checked by a proof assistant. The whole
development, from the definition of the household problem to the theorem that two equilibrium rates
of Aiyagari's log economy coincide, is written in Lean 4 \citep{Lean4} against the Mathlib library
\citep{Mathlib}, and the Lean kernel has verified that each proof follows from its hypotheses and
the standard axioms of the library. There are no unproved steps, no appeals to results that are
``well known'', and no numerical approximations. The formal statement of each result reported below
is named in the text, and Appendix~\ref{app:index} lists them with the files that contain them.

\paragraph{The main result.}
Fix Aiyagari's preferences and technology at $\mu=1$: log utility, $\beta=0.96$, Cobb--Douglas
production with capital share $\alpha=0.36$ and depreciation $\delta=0.08$, no borrowing. Let the
labour endowment follow any finite-state Markov chain whose transitions are monotone (a higher
state's row first-order stochastically dominates a lower state's) and whose lowest state is reached
from every state in one step. Then at most one interest rate in $(-\delta,\lambda)$, with
$\lambda=1/\beta-1$, is a stationary equilibrium rate (Theorem~\ref{thm:main}). Aiyagari's own
earnings processes, Tauchen discretisations of an AR(1) with positive persistence on any number of
grid points, satisfy the two conditions by construction (Corollary~\ref{cor:tauchen}). Outside
$(-\delta,\lambda)$ there is nothing: above $\lambda$ the household is patient and accumulates
without bound, so no stationary distribution has finite mean capital below demand
(Theorem~\ref{thm:above}); below $-4\%$ the only capital a log household holds is bounded by
$\beta\bar y/(1-\beta(1+r))$, which falls under demand as demand explodes at $-\delta$
(Theorem~\ref{thm:below}). Existence on any interval is reduced to a single inequality on capital
supply at its upper end (Theorem~\ref{thm:exist}), and the bisection Aiyagari ran is proved to
bracket and converge to the unique rate (Theorem~\ref{thm:bisect}).

\paragraph{Why the published arguments do not apply, and what does.}
The economics of the proof is the chain of \citet{Light2020}: individual saving rises with the
rate (his Theorem 1); therefore aggregate capital, the mean of the stationary distribution, rises
with the rate (his Theorem 2); therefore supply meets the strictly decreasing demand at most once
(his Theorem 3). Each link needed something new at $\beta=0.96$.

Theorem 1 is where relative risk aversion enters, through the monotonicity of $c\,u'(c)$ that the
rescaling argument needs, and for log utility that is available. But the argument is an induction
along value-function iterates, and it needs the borrowing constraint to be slack at the states it
compares and the asset cap, an artefact of compactness, to be slack everywhere. The published
sufficient conditions for both involve Lipschitz constants of the value function that scale like
$1/(1-\beta R)$, which at Aiyagari's rates is several hundred; at his parameters the tests fail by
orders of magnitude. What replaces them are two consequences of the Euler inequality itself: a
household at the constraint consumes at least the lowest income, which turns the corner condition
into the primitive inequality $\beta R\,u'(y_{\min})<u'(\text{cash on hand})$
(Section~\ref{sec:corner}); and a minimal marginal propensity to consume, $\kappa=1-\beta$ for log
utility at every rate, which bounds saving above and makes the cap slack under a single inequality
$2\beta\bar y<(1-\beta R)\bar a$ (Section~\ref{sec:mpc}).

Theorem 2 is a coupling: the two economies are run from a common initial distribution, the higher
rate's policy dominates pointwise, the Markov operator preserves the order, and the order passes to
the limit. The limit exists and is the stationary distribution because of a Doeblin condition that
the borrowing constraint supplies, an atom at zero assets after a run of bad shocks; the same
atom gives uniqueness of the stationary distribution at each rate, which is what turns capital
supply into a function of the rate (Section~\ref{sec:doeblin}).

Theorem 3 is elementary once supply is a monotone function, but Aiyagari's equilibrium condition is
stated at the equilibrium wage, and the household problem is solved at wage one. The bridge is
homogeneity of the stationary distribution in the wage, proved at the level of measures rather than
policies (Section~\ref{sec:homog}), so that the condition $K(r)=\E_r[a]$ at wage $w(r)$ becomes
$k(r)/w(r)=\E^{1}_r[a]$ for the unit-wage economy, and $k/w=\alpha/((1-\alpha)(r+\delta))$ needs no
root.

A last feature deserves mention because it shaped the whole development. Nothing is differentiated.
The value function of a constrained household is not known to be differentiable where the
constraint binds, and the standard envelope theorems do not apply there. The Euler inequalities
that drive every argument are derived from one-sided perturbations of the optimal plan and secant
slopes, in the manner of \citet{ClausenStrub2020}; Lipschitz bounds are secant bounds; monotonicity
is proved by comparing two plans, never by a derivative. This is not a stylistic choice: it is what
allowed the proofs to be completed at all in a setting where a proof assistant will not let a
differentiability assumption pass unnoticed.

\paragraph{Higher risk aversion.}
For $\mu\in\{3,5\}$ the argument fails at Theorem 1, and we show exactly how. The Euler-chain proof
of Theorem 1 reduces to a single \emph{elasticity condition} on the consumption function
(Section~\ref{sec:elasticity}): at tomorrow's assets $b$,
$R_1u'(c_1(b))\le R_2u'(c_2(b))$, that is, for CRRA, $(c_2/c_1)^\mu\le R_2/R_1$. For $\mu>1$ this
is false at high wealth, by the asymptotic linearity of the consumption function
\citep{MaToda2022}: the asymptotic marginal propensity to consume rises with the rate when $\mu>1$,
so consumption of the very rich rises with the rate by more than the condition allows. At
Aiyagari's numbers the condition holds for households with wealth up to about thirteen times mean
income at $\mu=3$ and seven times at $\mu=5$, and fails beyond. Saving of the rich nonetheless
rises with the rate; the sufficient condition simply does not see it.

The chain, it turns out, throws a term away. It starts from the supposition that saving falls with
the rate, which gives $c_2(a)>c_1(a)+(r_2-r_1)a$, and uses only $c_2(a)>c_1(a)$. Keeping the extra
interest income weakens the condition to $(c_2(b)/c_1(b))^\mu\le
(R_2/R_1)(1+(r_2-r_1)a/c_1(a))^\mu$ (Section~\ref{sec:slack}), and in the high-wealth limit this
holds for every $\mu$, exactly, with a margin in logarithmic terms of $1/(\kappa R)$ per unit of
rate difference, about $25$ at Aiyagari's parameters. The slack condition is therefore numerically
true at every wealth level. What we cannot do is prove it from primitives: it asks for the ratio of
consumption at two rates to a relative precision of the order of the rate difference, uniformly in
wealth, and no bound we have on the consumption function is that sharp. That is the frontier, and
we state it as such.

\paragraph{Huggett.}
In \citet{Huggett1993} the asset is a bond in zero net supply and households borrow. Against a
constant supply, monotone capital supply is not enough for single crossing, strict monotonicity is
needed, and we prove the strict form of Theorem 1 for creditors. For debtors a higher rate lowers
cash on hand, and saving can fall with the rate by up to the extra interest on the debt. We carry
this loss through the coupling and obtain $K_1\le K_2+(r_2-r_1)|\underline a|/(1-\Thorn_2)$, where
$\Thorn_2=(\beta R_2)^{1/\sigma}$ is the marginal propensity to save of the rich
(Section~\ref{sec:huggett}). At Huggett's numbers the amplification is about $1500$, and we have no
lower bound on the creditors' gain to set against it. We also show that the propensity-to-save
bound the argument needs cannot be obtained on a bounded asset region from the concavity of the
consumption function and its minimal-MPC floor, although it can on a ray, and we give the
counterexample as a theorem.

\paragraph{What is and is not claimed.}
The result is about Aiyagari's model as he wrote it, with one difference: assets are confined to
$[0,\bar a]$ for some cap $\bar a$, which his computation also imposes, and the theorem carries one
inequality saying the cap is large enough not to bind. Everything else is his: the utility, the
technology, the earnings process up to the two structural conditions. The uniqueness is of the
equilibrium interest rate; the stationary distribution at that rate is unique as well, by the
Doeblin argument. Existence is proved on any interval $[r_{lo},r_{hi}]\subset(-\delta,\lambda)$
whose upper end carries enough capital supply to exceed demand; that supply floor is the one
quantitative fact about the log household this paper does not prove, and Section~\ref{sec:exist}
says why. The bisection theorem takes an equilibrium as given and certifies the algorithm.

\paragraph{Related work.}
Uniqueness in Bewley economies is treated by \citet{Acikgoz2018}, \citet{Light2020},
\citet{Light2022}, and surveyed in \citet{TodaWalsh2024}; multiplicity at high risk aversion is
documented by \citet{YoungWalsh2025}. The consumption-function theory we lean on is
\citet{CarrollKimball1996} for concavity, \citet{MaToda2022} for asymptotic linearity and the
minimal MPC, \citet{CarrollShanker2026} for the buffer-stock bounds, and \citet{LiStachurski2014}
for the contraction in marginal-utility space that gives sufficiency of the Euler equation.
The formalisation uses the fixed-point
and monotone-operator results of \citet{SargentStachurski2024}, several of which we formalised
along the way (Section~\ref{sec:aux}). On formal verification in economics we know of no prior
formalisation of a recursive macroeconomic equilibrium; the closest work is the mean-field-games
literature on uniqueness, \citet{LasryLions2007}, \citet{GraberMatter2024},
\citet{CannarsaCapuani2018}, whose monotonicity conditions we discuss in Section~\ref{sec:aux}
and show do not reach stationary discounted equilibria.

Section~\ref{sec:model} sets out the model as formalised. Section~\ref{sec:main} states the
results for $\mu=1$. Section~\ref{sec:arch} explains the proof. Sections~\ref{sec:beyond} and
\ref{sec:huggett} give the frontier for $\mu>1$ and for Huggett. Section~\ref{sec:aux} collects
auxiliary results. Appendix~\ref{app:index} is the theorem index and Appendix~\ref{app:formal}
describes the formalisation.

\section{The model, as formalised}\label{sec:model}

\subsection{The household}

A household economy is a tuple
\[
  P=\bigl(Z,\;y,\;\Pi,\;r,\;\beta,\;u,\;\underline a,\;\bar a\bigr).
\]
$Z$ is a finite set of income states; $y:Z\to\R$ is the labour endowment with
$0<y_{\min}\le y(z)\le y_{\max}$; $\Pi$ is a stochastic matrix on $Z$; $r>-1$ is the interest rate;
$\beta\in[0,1)$ is the discount factor; $u$ is period utility; and $[\underline a,\bar a]$ with
$\underline a\le\bar a$, $0\le \bar a$, is the asset region. The borrowing limit and the cap are
\emph{type parameters} of the formal object, so that Aiyagari's economy ($\underline a=0$) and
Huggett's ($\underline a<0$) are the same structure, and a theorem proved for a general floor
applies to both. The state is $(a,z)\in[\underline a,\bar a]\times Z$, cash on hand is
$m(a,z)=y(z)+(1+r)a$, and the household solves
\begin{equation}\label{eq:bellman}
  V(a,z)=\max_{a'\in[\underline a,\;\min\{\bar a,\,m(a,z)\}]}
  \Bigl\{u\bigl(m(a,z)-a'\bigr)+\beta\sum_{z'}\Pi(z,z')V(a',z')\Bigr\}.
\end{equation}
Utility is required to be continuous, increasing and strictly concave on a domain that is either
$(0,\infty)$ (utility unbounded below, as for log and CRRA with $\mu\ge1$) or $[0,\infty)$; off the
domain the reward is $-\infty$. With a negative borrowing limit the model carries Aiyagari's
natural-borrowing-limit condition, $y_{\min}+r\underline a>0$: income at the limit covers the
interest on the debt. Because raising $r$ raises the debt service, a rate is \emph{admissible}
for $P$, written $\operatorname{RateOK}(r)$, when $1+r>0$ and the limit remains sustainable at $r$;
with $\underline a=0$ the second condition is vacuous. $P.\text{withRate}(r)$ denotes the same
economy at rate $r$.

The value function exists and is the unique bounded fixed point of \eqref{eq:bellman}; the policy
$a'=g(a,z)$ is unique by strict concavity; consumption is $c(a,z)=m(a,z)-g(a,z)$. These are
theorems of the development (a stochastic principle of optimality with extended-real rewards), not
assumptions. Where a result needs consumption to be positive at the optimum it takes the hypothesis
$\operatorname{PositiveConsumption}$, which is discharged automatically when utility is unbounded
below and by a marginal Inada condition otherwise.

\subsection{Distributions}

The state space $S=[\underline a,\bar a]\times Z$ is compact. The policy and the transition matrix
define a Markov kernel on $S$, and its action on Borel probability measures is the operator
$\mu\mapsto T\mu$. A \emph{stationary distribution} is a fixed point, $T\mu=\mu$. Aggregate capital
is mean assets,
\[
  K(\mu)=\int a\,d\mu(a,z).
\]
Existence of a stationary distribution is by the Krylov--Bogolyubov argument on the compact state
space, a theorem of the development; uniqueness is not automatic and is discussed in
Section~\ref{sec:doeblin}.

Two structural conditions on the earnings process recur. The transitions are \emph{monotone} when
for any $z_1,z_2$ with $y(z_1)\le y(z_2)$, the row $\Pi(z_2,\cdot)$ first-order stochastically
dominates $\Pi(z_1,\cdot)$ in the income order; iid income is the trivial case, and a Tauchen
discretisation of an AR(1) with $\rho\ge0$ is monotone by construction. A state $z_0$ is
\emph{reached from everywhere} when $\Pi(z,z_0)>0$ for every $z$; for the Doeblin argument $z_0$
is a state of minimal income.

\subsection{The firm and the equilibrium condition}

Aiyagari's firm is Cobb--Douglas with capital share $\alpha$ and depreciation $\delta$, with labour
supplied inelastically and the endowment process normalised to mean one, so that per-capita capital
demand and the wage at rate $r$ are
\[
  k(r)=\Bigl(\frac{\alpha}{r+\delta}\Bigr)^{1/(1-\alpha)},\qquad
  w(r)=(1-\alpha)k(r)^{\alpha},\qquad
  \frac{k(r)}{w(r)}=\frac{\alpha}{(1-\alpha)(r+\delta)}.
\]
His stationary equilibrium condition is $k(r)=\E_{r,w(r)}[a]$: capital demand equals mean asset
holdings under the stationary distribution of the household economy at rate $r$ and wage $w(r)$.
The household problem at wage $w$ is the problem at wage one with incomes and the cap scaled by
$w$, and the stationary distribution scales with it, so the condition is equivalent to
\begin{equation}\label{eq:eqm}
  \E^{1}_{r}[a]=D(r),\qquad D(r):=\frac{k(r)}{w(r)}=\frac{\alpha}{(1-\alpha)(r+\delta)},
\end{equation}
where $\E^1_r$ is taken under a stationary distribution of the unit-wage economy at rate $r$. The
formal object $\operatorname{normalisedDemand}(\alpha,\delta,r)$ is $D(r)$, and an
\emph{equilibrium rate} is an admissible $r$ together with a stationary distribution $\mu$ of
$P.\text{withRate}(r)$ satisfying $K(\mu)=D(r)$. The reduction from Aiyagari's form to
\eqref{eq:eqm} is Section~\ref{sec:homog}.

$D$ is strictly decreasing on any interval with $r+\delta>0$ and unbounded as $r\downarrow-\delta$.
The rate of time preference is $\lambda=1/\beta-1$; at $\beta=0.96$ it is $1/24\approx4.17\%$, and
Aiyagari's reported equilibrium rates at $\mu=1$ (his Table II) lie between $3.31\%$ and $4.17\%$; at $\mu\in\{3,5\}$ they range down to $-0.35\%$.

\subsection{Aiyagari's calibration}

Throughout, ``Aiyagari's numbers'' means $\beta=24/25$, $\alpha=9/25$, $\delta=2/25$,
$\underline a=0$, log utility for $\mu=1$ and $u(c)=c^{1-\mu}/(1-\mu)$ for $\mu\in\{3,5\}$, and a
labour endowment $\log l'=\rho\log l+\sigma\sqrt{1-\rho^2}\,\varepsilon$ with
$\rho\in\{0,0.3,0.6,0.9\}$, $\sigma\in\{0.2,0.4\}$, discretised by Tauchen's method on
$[-3\sigma,3\sigma]$. The formal statements use the rational values; $0.96$ is $24/25$ exactly.

\section{Results for the log calibration}\label{sec:main}

Throughout this section $P$ is a household economy with $\underline a=0$, cap $\bar a$, log
utility, $\beta=24/25$, and an earnings process $(Z,y,\Pi)$. Two conditions on the earnings
process are assumed in every uniqueness statement:
\begin{enumerate}[label=(E\arabic*)]
  \item\label{E1} the transitions are monotone;
  \item\label{E2} there is a state $z_0$ with $y(z_0)=y_{\min}$ and $\Pi(z,z_0)>0$ for every $z$.
\end{enumerate}
The cap enters through one inequality, evaluated at the larger of the two rates compared:
\begin{equation}\label{eq:cap}
  2\beta\,y_{\max}<\bigl(1-\beta(1+r)\bigr)\,\bar a .
\end{equation}
It says the cap was set large enough not to bind; nothing else about $\bar a$ is used.

\subsection{Uniqueness}

\begin{theorem}[Aiyagari at $\mu=1$: the equilibrium rate is unique on $(-\delta,\lambda)$]
\label{thm:main}
Let $r_1,r_2\in(-2/25,\,1/24)$ be admissible, and suppose that for $i=1,2$ there is a stationary
distribution $\mu_i$ of $P.\text{withRate}(r_i)$ with $K(\mu_i)=D(r_i)$, where
$D(r)=\alpha/((1-\alpha)(r+\delta))$ with $\alpha=9/25$, $\delta=2/25$. If \eqref{eq:cap} holds at
$\max\{r_1,r_2\}$, then $r_1=r_2$.
\formal{IncomeFluctuation.aiyagari1994\_log\_equilibriumRate\_unique\_Ioo}
\end{theorem}

The theorem is stated for every interval $[r_{lo},r_{hi}]$ with $-\delta<r_{lo}\le r_{hi}<\lambda$
in the form \lean{aiyagari1994\_log\_equilibriumRate\_unique}, with \eqref{eq:cap} at $r_{hi}$;
the open-interval version applies it to $[\min\{r_1,r_2\},\max\{r_1,r_2\}]$. The lower end
$-\delta$ is where demand becomes infinite; the upper end $\lambda$ is where the household becomes
patient. Both are excluded for a reason (Theorems~\ref{thm:above} and \ref{thm:below}), so
Theorem~\ref{thm:main} covers every rate at which an equilibrium could exist, up to a band above
$\lambda$ discussed below.

\begin{corollary}[Aiyagari's own earnings processes]\label{cor:tauchen}
Let $n\ge0$, $\rho\in[0,1)$, $\sigma>0$, $c>0$, and let the endowment take the $n+2$ values
$c\,e^{-3\sigma+i\cdot 6\sigma/(n+1)}$, $i=0,\dots,n+1$, with the Tauchen transition matrix of
$\log l'=\rho\log l+\sigma\sqrt{1-\rho^2}\,\varepsilon$ on that grid. Then the conclusion of
Theorem~\ref{thm:main} holds, with \eqref{eq:cap} read with $y_{\max}=c\,e^{3\sigma}$.
\formal{aiyagariTauchen\_equilibriumRate\_unique\_Ioo}
\end{corollary}

The corollary covers all eight of Aiyagari's processes and any finer grid. Its proof is that the
Tauchen chain satisfies \ref{E1} and \ref{E2}: the cumulative rows
$F_i(k)=\Phi\bigl((x_k+h/2-\rho x_i)/\sigma_\varepsilon\bigr)$ are decreasing in $i$ because $\Phi$
is monotone and $\rho\ge0$, which is first-order dominance, and every entry of the first column is
a value of $\Phi$ on the real line, hence positive. The normal distribution is Mathlib's
$\mathrm{gaussianReal}$, so $\Phi$ is a genuine cumulative distribution function and the
positivity is a theorem about the Lebesgue integral, not a numerical fact. The scale $c$ is
whatever normalises the mean to one; uniqueness does not depend on it except through
\eqref{eq:cap}.

\begin{remark}[What is used about the earnings process]
Only \ref{E1} and \ref{E2}. No transition probability is named, no stationary distribution of the
chain is computed, and the number of states is arbitrary. \ref{E1} is what Light's Theorem 1 needs
to propagate monotonicity of the value function in income through the expectation; \ref{E2} is
what the Doeblin argument needs to produce an atom. For $\rho=0$ both are immediate; for $\rho>0$
they are Huggett's Lemma 1 for Tauchen chains.
\end{remark}

\subsection{No equilibrium outside \texorpdfstring{$(-\delta,\lambda)$}{the interval}}

\begin{theorem}[Above the rate of time preference]\label{thm:above}
Let $r\ge0$ with $\beta(1+r)>1$, and let $\mu$ be stationary for $P.\text{withRate}(r)$. If
\[
  y_{\max}+(1+r)D<\Bigl(1-\frac{1}{\beta(1+r)}\Bigr)\bigl(y_{\min}+r\bar a\bigr),
\]
then $K(\mu)\ne D$. In particular no equilibrium exists at $r$ once $r-\lambda$ exceeds a
quantity of order $(y_{\max}+(1+r)D(r))/(\beta r\bar a)$.
\formal{IncomeFluctuation.log\_no\_equilibrium\_of\_patient\_risk}
\end{theorem}

The proof is a supermartingale bound. With $\beta R>1$ the Euler inequality makes
$u'(c_t)$ a supermartingale that cannot decrease on average, so consumption cannot rise on average,
while cash on hand grows at rate $r$ on the part of the distribution near the cap; integrating
against the stationary distribution gives a floor on $K(\mu)$ of order $(1-1/\beta R)\,r\bar a$.
The band $(\lambda,\lambda+\epsilon)$ that escapes has width shrinking like $1/\bar a$, and with
$\bar a$ chosen to satisfy \eqref{eq:cap} it is thin. A version for CRRA with integer risk
aversion $n$, via Jensen on $x^{-n}$, is \lean{crra\_no\_equilibrium\_of\_patient\_mean}.

\begin{theorem}[Negative rates]\label{thm:below}
Assume \ref{E2} and that the earnings process is normalised so that the invariant distribution of
$\Pi$ has mean at most one. Let $\mu$ be stationary for $P.\text{withRate}(r)$.
\begin{enumerate}[label=(\roman*)]
\item For any period utility satisfying the standing assumptions, if $r\in(-2/25,-13/250]$ then
  $K(\mu)\ne D(r)$.
\item For log utility, if $24\,y_{\max}<\bar a$ and $r\in(-2/25,-1/25]$ then $K(\mu)\ne D(r)$.
\end{enumerate}
\formal{IncomeFluctuation.aiyagari1994\_no\_equilibrium\_of\_budget; IncomeFluctuation.aiyagari1994\_log\_no\_equilibrium\_below\_mean}
\end{theorem}

Part (i) uses nothing about preferences: at a negative rate, integrating the budget constraint
against a stationary distribution gives $K\le \bar Y/(-r)$, the capital whose interest cost
exhausts mean income, which is the zero-consumption supply schedule of \citet{YoungWalsh2025}.
Part (ii) uses the log ceiling $K\le\beta\bar Y/(1-\beta(1+r))$ from the minimal MPC
(Section~\ref{sec:mpc}). Both need the income marginal of the stationary distribution to be the
chain's invariant distribution, which is where \ref{E2} enters, through a finite-chain uniqueness
argument (\lean{invariant\_eq\_of\_reach}). In both cases the ceiling falls below $D(r)$, which
grows without bound as $r\downarrow-\delta$.

Together, Theorems~\ref{thm:main}--\ref{thm:below}: for Aiyagari's log calibration the set of
equilibrium rates is contained in $(-1/25,\lambda+\epsilon)$ and has at most one point in
$(-\delta,\lambda)$.

\subsection{Existence}\label{sec:exist}

\begin{theorem}[Existence on an interval, given a supply floor at its top]\label{thm:exist}
Let $-\delta<r_{lo}<r_{hi}<\lambda$, assume \ref{E1}, \ref{E2} and \eqref{eq:cap} at $r_{hi}$.
Suppose
\[
  \frac{\beta\,y_{\max}}{1-\beta(1+r_{lo})}\le D(r_{lo}),\qquad
  D(r_{hi})\le m\le K(\mu)\ \text{ for every stationary }\mu\text{ at }r_{hi}.
\]
Then there is an equilibrium rate in $[r_{lo},r_{hi}]$.
\formal{IncomeFluctuation.aiyagari1994\_log\_exists\_equilibrium\_of\_floor}
\end{theorem}

The first inequality is the log capital ceiling at $r_{lo}$ against demand there, and it holds
for $r_{lo}$ near $-\delta$ because demand is unbounded. The second is a floor $m$ on capital
supply at $r_{hi}$ that exceeds demand there, and it is the one quantitative fact about the log
household that this development does not prove. Aiyagari's equilibria sit where $k/w\approx5$; the
analytic supply floors available (an atom at the cap, positive saving from the corner) are nowhere
near that, and how much an impatient household accumulates before impatience takes over is a
quantitative question about the policy that the qualitative theory does not answer. With the floor
supplied, existence follows by the intermediate value theorem, because capital supply along the
unique stationary distribution is continuous in the rate (Section~\ref{sec:cont}).

\subsection{Computation}

\begin{theorem}[Bisection converges to the equilibrium rate]\label{thm:bisect}
Under the hypotheses of Theorem~\ref{thm:main} on $[r_{lo},r_{hi}]$, let $\nu$ be any selection of
stationary distributions, $\nu(r)$ stationary for $P.\text{withRate}(r)$ on $[r_{lo},r_{hi}]$,
and let $E(r)=K(\nu(r))-D(r)$ be excess supply. If $r^*\in[r_{lo},r_{hi}]$ is an equilibrium rate,
then the bisection on $E$ started from $[r_{lo},r_{hi}]$ brackets $r^*$ at every step, its
midpoint after $n$ steps is within $(r_{hi}-r_{lo})/2^{n+1}$ of $r^*$, and its endpoints converge
to $r^*$.
\formal{IncomeFluctuation.aiyagari1994\_log\_bisection}
\end{theorem}

Excess supply is strictly increasing under these hypotheses, because supply is monotone
(Section~\ref{sec:coupling}) and demand strictly decreasing, and strict monotonicity alone decides
which half of a bracket contains the zero; neither continuity nor a sign condition at the ends is
used. Any stationary distribution at $r$ is the selection's, by the Doeblin uniqueness, so the
algorithm does not depend on how the stationary distribution is computed. The theorem certifies
the outer loop of the standard procedure, as in \citet{KirkbyToolkit2017}, given exact inner
solutions; convergence of the discretised inner solutions to the true ones is the subject of
\citet{Kirkby2017} and \citet{Kirkby2019}, and is not re-proved here.

\section{How the proof goes}\label{sec:arch}

The argument is Light's chain. This section describes each link and what had to be done to make
it hold at $\beta=0.96$. Formal names are given at each step; the emphasis is on the economics.

\subsection{Theorem 1: saving rises with the rate}\label{sec:thm1}

Fix two admissible rates $r_1<r_2$ and write $g_i$ for the policy at $r_i$. Light's Theorem 1 is
$g_1(a,z)\le g_2(a,z)$ for every state. His proof compares the two Bellman operators along
value-function iterates: the map $v\mapsto v(ta,\cdot)$, $t=R_1/R_2$, turns a household at the
higher rate with assets $ta$ into one with the same cash on hand as a household at the lower rate
with assets $a$, and the induction shows that the difference of the two value functions has
increasing differences in assets, which is what makes the higher-rate policy larger. The step where
the utility function enters is an inequality of the form $t\,u'(c(ta))\le u'(c(a))$ for the
rescaled consumption, and this is where relative risk aversion at most one is used: it is
monotonicity of $c\,u'(c)$, which for log utility is an identity.
(\lean{monotoneOn\_bellman\_sub\_scaled\_of\_piv}, \lean{policy\_mono\_withRate\_crra}.)

The induction needs the constraint slack at the compared states and the cap slack everywhere along
the iterates, and the published tests for both scale like $1/(1-\beta R)$. At Aiyagari's rates
$1-\beta R$ is about $0.002$, and the tests fail by orders of magnitude. What replaces them is in
the next two subsections. With those in hand the induction is run locally, on an interval of rates
short enough that the constants survive, and the local statements are chained across
$[r_{lo},r_{hi}]$ (\lean{LargeBetaLight}).

\subsection{The corner from the Euler inequality}\label{sec:corner}

\begin{lemma}[Nobody at the constraint consumes less than the lowest income]
At an interior choice $g(a,z)>0$ the Euler inequality reads $u'(c)\le\beta R\,\E[u'(c')]$. Take the
income state at which consumption at zero assets is smallest. If that household saved, tomorrow's
consumption would be at least today's minimum, so $u'(c)\le\beta R\,u'(c)$, impossible with
$\beta R<1$. So it is at the corner, consumes its income, and consumption at the constraint is at
least $y_{\min}$ in every state.
\formal{IncomeFluctuation.minIncome\_le\_consumptionFn\_floor}
\end{lemma}

The corner condition at a state $(a,z)$ is then the contrapositive of the Euler inequality: if
$\beta R\,u'(y_{\min})<u'(m(a,z))$, saving is not optimal, because saving would give
$u'(m(a,z))\le u'(c)\le\beta R\,\E[u'(c')]\le\beta R\,u'(y_{\min})$. This is a primitive
inequality, checkable by hand at $\beta=0.96$, and it is what
\lean{aiyagari1994\_log\_equilibriumRate\_unique} evaluates internally with an explicit corner
level $a_0=\min\{\bar a,\;y_{\min}(1-\beta R)/(2\beta R^2)\}$.
(\lean{EulerCorner}.)

\subsection{The cap from the minimal marginal propensity to consume}\label{sec:mpc}

\begin{lemma}[Minimal MPC]
For CRRA utility with parameter $\gamma$ at a zero borrowing limit, with $\beta R<1$,
\[
  c(a,z)\ \ge\ \kappa\, m(a,z),\qquad \kappa=1-\frac{\Thorn}{R},\quad \Thorn=(\beta R)^{1/\gamma}.
\]
\formal{IncomeFluctuation.crra\_minMPC\_mul\_le\_consumptionFn}
\end{lemma}

$\Thorn$ is Carroll's absolute patience factor and $\kappa$ the MPC of a perfect-foresight consumer
with no human wealth \citep{CarrollShanker2026}. The bound is self-reproducing under the Bellman
operator: if the continuation's consumption satisfies it, then at an interior choice the Euler
inequality gives $c^{-\gamma}\le\beta R(\kappa R a')^{-\gamma}$, hence $c\ge(\kappa R/\Thorn)a'$, and
$(\kappa R/\Thorn)/(1+\kappa R/\Thorn)=\kappa$ exactly because $1-\kappa=\Thorn/R$; at the corner
the household consumes everything. For log utility $\kappa=1-\beta$ at every rate, which is why the
log results carry no rate-dependent constant.

Two things follow. Saving is at most $(1-\kappa)m(a,z)$, so at the cap
$g(\bar a,z)\le\beta(y_{\max}+R\bar a)<\bar a$ whenever \eqref{eq:cap} holds, and the cap is slack
along every iterate and at the fixed point (\lean{crra\_policy\_lt\_cap\_of\_minMPC}). And
integrating against a stationary distribution, $K\le\beta(\bar Y+RK)$, so
\begin{equation}\label{eq:ceiling}
  K(\mu)\le\frac{\beta\,\bar Y}{1-\beta R},
\end{equation}
the log capital ceiling used in Theorems~\ref{thm:below} and \ref{thm:exist}
(\lean{log\_aggregateCapital\_le}, \lean{log\_aggregateCapital\_le\_mean}).

\subsection{Theorem 2: the coupling}\label{sec:coupling}

\begin{proposition}[Aggregate capital rises with the rate]
Let $g_1\le g_2$ pointwise on the region, and let the two Markov operators be run from a common
initial distribution $\mu_0$, converging to stationary distributions $\mu_1$ and $\mu_2$. Then
$\mu_2$ dominates $\mu_1$ in the asset coordinate, and $K(\mu_1)\le K(\mu_2)$.
\formal{IncomeFluctuation.monotoneOn\_capitalSupply\_of\_policy\_mono}
\end{proposition}

Dominance is tested against bounded continuous functions that are increasing in assets at each
income state. Three facts carry the argument: the Markov operator preserves that test class, because
the policy is increasing in assets (\lean{policy\_mono}); a pointwise-larger policy gives a
dominating operator on that class; and dominance is closed under weak limits. The asset coordinate
itself is in the test class, so the mean is ordered. This is the argument of \citet{Light2022},
Theorem 2, which asks for monotonicity in one coordinate only; monotonicity in the whole state
fails here because the policy need not rise with the income shock.

\subsection{Uniqueness and convergence of the stationary distribution}\label{sec:doeblin}

The coupling needs the iterates from $\mu_0$ to converge, and the equilibrium condition needs
capital supply to be a function of the rate, not a correspondence. Both come from a Doeblin
condition. The asset transition is deterministic given the shock, so two households with different
assets never meet, and no minorisation against a spread-out measure holds. The borrowing constraint
supplies an atom: if a run of $N$ draws of the worst income state $z_0$ drives every household to
exactly zero assets, then every $N$-step transition charges the single point $(0,z_0)$ with
probability at least $\Pi_{\min}(\cdot,z_0)^N$, which is a Doeblin condition. The contraction is run
on the function side: the oscillation of $T^N h$ is at most $(1-\epsilon)$ times that of $h$, so
two stationary distributions agree on every bounded continuous function, and the iterates from any
start converge weakly to the stationary distribution.
(\lean{stationary\_unique}, \lean{tendsto\_pushProb\_iterate}.)

That a run of bad shocks reaches the constraint is the \emph{decline} property: above some asset
level the household dissaves in the worst state. For an impatient household this is
\citet{Acikgoz2018}, Proposition 4, and it is proved here from the corner condition and the minimal
MPC (\lean{ImpatientDecline}). \ref{E2} is what makes the run have positive probability from every
state. At each rate in $[r_{lo},r_{hi}]$ the conclusion is a unique stationary distribution,
reached from every start (\lean{log\_existsUnique\_isStationary}).

\subsection{Theorem 3: single crossing, and Aiyagari's form of the condition}\label{sec:homog}

With supply a monotone function of the rate on $[r_{lo},r_{hi}]$ and $D$ strictly decreasing, two
equilibrium rates coincide: this is \lean{eq\_of\_monotoneOn\_of\_strictAntiOn} and needs no
continuity. The remaining step is to connect \eqref{eq:eqm} to Aiyagari's condition at the
equilibrium wage.

\begin{proposition}[Homogeneity of the stationary distribution]
Scaling every income and the cap by $w>0$ scales the policy by $w$. The map $(a,z)\mapsto(wa,z)$
therefore commutes with the Markov operator, so it carries stationary distributions of the unit-wage
economy to stationary distributions of the wage-$w$ economy, and aggregate capital scales by $w$.
\formal{IncomeFluctuation.HomotheticDistribution}
\end{proposition}

This is Açıkgöz's Proposition 7 at the level of measures. It is what lets the theorems be stated in
Aiyagari's own form, $K(r)=\E_r[a]$ at $w(r)$, while the household side is solved once at wage one
and the firm side reduces to $k/w=\alpha/((1-\alpha)(r+\delta))$, which needs no root
(\lean{aiyagari1994\_equilibriumRate\_unique\_wage}).

\subsection{Continuity of supply, for existence}\label{sec:cont}

Existence uses the intermediate value theorem, which needs capital supply continuous in the rate.
The rate is not a parameter of the household problem in the formal development but part of its
state, so that a family of problems indexed by the rate is one problem on a product space, and
continuity of the policy in the rate is joint continuity of one policy. Weak continuity of the
stationary distribution in the rate then follows from the Doeblin uniqueness (a limit of stationary
distributions is stationary, and there is only one), and aggregation is weakly continuous
because the asset coordinate is bounded and continuous (\lean{continuousOn\_capitalSupply}).

\subsection{What is never differentiated}\label{sec:deriv}

Every inequality above is derived without a derivative of the value or policy function. The Euler
inequalities come from perturbing the optimal plan by $\epsilon$ in one direction and letting
$\epsilon$ shrink: carrying the poorer household's plan forward, or the richer household's plan
back, gives one-sided bounds on the marginal value of assets, and the first-order condition is read
off the pair (\lean{euler\_le}, \lean{euler\_ge}; the construction is that of
\citealp{ClausenStrub2020}). Lipschitz constants are secant bounds; monotonicity in the rate is a
comparison of two plans. Where the development does establish differentiability of the value
function (\lean{IncomeFluctuationEnvelope}) it uses the Clausen--Strub sandwich and not the
Benveniste--Scheinkman argument, and the uniqueness proof does not depend on it.

\section{Higher risk aversion: the frontier for \texorpdfstring{$\mu\in\{3,5\}$}{mu = 3, 5}}
\label{sec:beyond}

Everything in Section~\ref{sec:arch} except Theorem 1 goes through for CRRA utility with any
$\mu>0$: the corner, the minimal MPC with $\kappa=1-\Thorn/R$, the cap slack, the Doeblin
uniqueness and convergence, the coupling, the homogeneity, single crossing. The formal statement is
\lean{crra\_equilibriumRate\_unique\_of\_marginal}, and it carries exactly one hypothesis that is
not discharged: Theorem 1. This section says what that hypothesis is, why it fails at Aiyagari's
numbers, what it becomes when the proof is tightened, and where that leaves the problem.

\subsection{Theorem 1 reduces to an elasticity condition}\label{sec:elasticity}

\begin{proposition}[Saving rises with the rate under an elasticity condition]\label{prop:elas}
Let $r_1\le r_2$ be admissible, with positive consumption and a slack cap at both rates, and let
$u'$ be strictly decreasing. Fix $(a,z)$ with $a\ge0$ and let $b=g_1(a,z)$. If
\begin{equation}\label{eq:EC}
  R_1\,u'\bigl(c_1(b,z')\bigr)\ \le\ R_2\,u'\bigl(c_2(b,z')\bigr)\qquad\text{for every }z',
\end{equation}
then $g_1(a,z)\le g_2(a,z)$.
\formal{IncomeFluctuation.policy\_le\_policy\_withRate\_of\_marginal\_at}
\end{proposition}

The proof is four lines of Euler inequalities and holds for any period utility. Suppose
$g_2(a,z)<g_1(a,z)=b$. Then the higher-rate household has room under the cap, so
$\beta R_2\,\E[u'(c_2(g_2,\cdot))]\le u'(c_2(a,z))$; the lower-rate household has room above the
floor, so $u'(c_1(a,z))\le\beta R_1\,\E[u'(c_1(b,\cdot))]$; it consumes strictly less today, so
$u'(c_2(a,z))<u'(c_1(a,z))$; and consumption rises with assets, so
$u'(c_2(b,\cdot))\le u'(c_2(g_2,\cdot))$. Chaining, $R_2\E[u'(c_2(b,\cdot))]<R_1\E[u'(c_1(b,\cdot))]$,
which \eqref{eq:EC} forbids. For CRRA, \eqref{eq:EC} is
\[
  \Bigl(\frac{c_2(b,z')}{c_1(b,z')}\Bigr)^{\mu}\le\frac{R_2}{R_1},
\]
an elasticity of consumption to the gross rate of at most $1/\mu$ (\lean{crra\_elasticity\_iff}).
It says that the income effect of the rate on consumption does not outweigh the human-wealth
effect. Uniqueness on $[r_{lo},r_{hi}]$ follows from \eqref{eq:EC} at every state and every pair of
rates (\lean{crra\_equilibriumRate\_unique\_of\_marginal}), and, since the coupling only visits
states the economy visits, from \eqref{eq:EC} on any interval $[0,\bar a']$ that the lower-rate
policy maps into itself (\lean{crra\_equilibriumRate\_unique\_of\_marginal\_on},
Section~\ref{sec:invariant}).

\subsection{The condition fails for the rich when \texorpdfstring{$\mu>1$}{mu > 1}}

By \citet{MaToda2022}, the consumption function is asymptotically linear, $c(x)/x\to\kappa(R)$,
with $\kappa(R)=1-(\beta R)^{1/\mu}/R$ the minimal MPC. For $\mu>1$, $\kappa$ rises with $R$. So at
high wealth $c_2(b)/c_1(b)\to\kappa_2R_2/(\kappa_1R_1)$ and
\[
  \Bigl(\frac{\kappa_2R_2}{\kappa_1R_1}\Bigr)^{\mu}
  =\Bigl(\frac{\kappa_2}{\kappa_1}\Bigr)^{\mu}\Bigl(\frac{R_2}{R_1}\Bigr)^{\mu}>\frac{R_2}{R_1},
\]
so \eqref{eq:EC} fails for all large $b$ (\lean{crra\_elasticity\_fails\_of\_asymptotic\_mpc}).
Saving of the rich nonetheless rises with the rate, since $g\approx\Thorn b$ with
$\Thorn=(\beta R)^{1/\mu}$ increasing in $R$; the sufficient condition does not see it. With the
linear approximation $c\approx\kappa(R)(Rb+H)$ and $H\approx\bar Y/r$, the human-wealth effect
dominates and consumption falls with the rate for $b$ below about $13\bar Y$ at $\mu=3$ and
$7\bar Y$ at $\mu=5$, where \eqref{eq:EC} holds; above, it fails. The stationary distribution
reaches far beyond those levels, because the top income type keeps accumulating until wealth of
order $y_{\max}/(1-\Thorn)$, so the invariant-interval form does not rescue the argument at
Aiyagari's numbers.

\subsection{The chain throws a term away}\label{sec:slack}

The supposition $g_2(a,z)<g_1(a,z)$ gives not $c_2(a,z)>c_1(a,z)$ but
\[
  c_2(a,z)>c_1(a,z)+(r_2-r_1)\,a,
\]
and the proof above uses only the weaker inequality. Keeping the extra interest income, the chain
closes under a weaker condition.

\begin{proposition}[The slack elasticity condition]\label{prop:slack}
In the setting of Proposition~\ref{prop:elas}, with $u'>0$, if
\begin{equation}\label{eq:SEC}
  R_1\,u'\bigl(c_1(b,z')\bigr)\,u'\bigl(c_1(a,z)+(r_2-r_1)a\bigr)
  \ \le\ R_2\,u'\bigl(c_2(b,z')\bigr)\,u'\bigl(c_1(a,z)\bigr)\qquad\text{for every }z',
\end{equation}
then $g_1(a,z)\le g_2(a,z)$. For CRRA, \eqref{eq:SEC} reads
\[
  \Bigl(\frac{c_2(b,z')}{c_1(b,z')}\Bigr)^{\mu}\le\frac{R_2}{R_1}
  \Bigl(1+\frac{(r_2-r_1)\,a}{c_1(a,z)}\Bigr)^{\mu}.
\]
\formal{IncomeFluctuation.policy\_le\_policy\_withRate\_of\_slack\_at; crra\_slack\_elasticity\_iff}
\end{proposition}

The new factor is at least one and grows with $a/c_1(a)$, which is largest exactly where
\eqref{eq:EC} fails. In the high-wealth limit, with $L_i=\kappa_iR_i=R_i-\Thorn_i$ the asymptotic
slopes and $b\approx\Thorn_1a$, $c_1(a)\approx L_1a$, both sides have limits and the inequality
between them is exact:

\begin{proposition}[The slack condition holds in the homogeneous limit for every $\mu$]
\label{prop:homog}
For $0<R_1\le R_2$, $\Thorn_1\le\Thorn_2$, $\Thorn_1<R_1$, $\Thorn_2<R_2$ and $\mu\ge0$,
\[
  \Bigl(\frac{R_2-\Thorn_2}{R_1-\Thorn_1}\Bigr)^{\mu}\ \le\
  \frac{R_2}{R_1}\Bigl(\frac{R_2-\Thorn_1}{R_1-\Thorn_1}\Bigr)^{\mu},
\]
with strict inequality when $\Thorn_1<\Thorn_2$ and $\mu>0$. Consequently, along tomorrow's
assets, if $c_2(b(a))/c_1(b(a))\to L_2/L_1$ and $c_1(a)/a\to L_1$, the CRRA slack condition holds
for all large $a$.
\formal{crra\_slack\_elasticity\_of\_patience; crra\_slack\_elasticity\_holds\_of\_asymptotic\_mpc}
\end{proposition}

The proof is one line: $R_2-\Thorn_2\le R_2-\Thorn_1$ because $\Thorn_1\le\Thorn_2$, and
$R_2/R_1\ge1$ is spare. In logarithmic terms per unit of rate difference the two sides grow at
rates $\mu(\kappa'R+\kappa)/(\kappa R)$ and $1/R+\mu/(\kappa R)$, whose difference is exactly
$1/(\kappa R)$ for every $\mu$: about $25$ at Aiyagari's numbers. Table~\ref{tab:slack} gives the
values.

\begin{table}[t]
\centering
\caption{The slack elasticity condition at Aiyagari's numbers, $r\approx4\%$. Rows two and three
are the logarithmic growth rates of the two sides of \eqref{eq:SEC} per unit of rate difference in
the high-wealth limit; the condition holds when the left is below the right. Row four is the wealth
below which consumption falls with the rate, where the unweakened condition \eqref{eq:EC} already
holds.}
\label{tab:slack}
\begin{tabular}{lcc}
\toprule
 & $\mu=3$ & $\mu=5$\\
\midrule
minimal MPC $\kappa$ & $0.0390$ & $0.0388$\\
left side, deep tail & $50$ & $100$\\
right side, deep tail & $75$ & $125$\\
wealth below which \eqref{eq:EC} holds (multiples of mean income) & $\approx13$ & $\approx7$\\
\bottomrule
\end{tabular}
\end{table}

So at every wealth level the slack condition is, numerically, true: in the body because
consumption falls with the rate, in the tail by Proposition~\ref{prop:homog} with margin, and in
between with margin as well. The reduction theorems with the slack condition in place of
\eqref{eq:EC} are \lean{equilibriumRate\_unique\_of\_slack},
\lean{equilibriumRate\_unique\_of\_slack\_on} and
\lean{crra\_equilibriumRate\_unique\_of\_slack\_on}.

\subsection{What remains}\label{sec:frontier}

To prove \eqref{eq:SEC} from primitives one needs the ratio $c_2(b)/c_1(b)$ of consumption at two
rates to a relative precision of the order of $r_2-r_1$, uniformly in $b$ on the support, since
that is the size of the margin. Two available routes do not deliver it. The sandwich
$\kappa x\le c(x)\le\kappa x+C$, with $C$ of the order of mean income, has a gap of order one and
so pins the ratio only for rate pairs far enough apart, whereas uniqueness must handle all pairs.
The fixed-point perturbation bound for the Coleman operator in the marginal-utility metric of
\citet{LiStachurski2014}, which we formalised (Section~\ref{sec:aux}), gives a bound that is
uniform in wealth in the wrong sense: it loses a factor of the inverse contraction modulus,
$1/(1-\Thorn)\approx2000$ at these numbers, and is vacuous where $u'$ is small. What would be
needed is a Lipschitz estimate in the rate for the consumption function that is relative rather than
absolute, that is, comparative statics of $c$ in $r$ uniform on the support. We know of no such
result. The status of $\mu\in\{3,5\}$ is therefore: reduced to an inequality on two consumption
functions that is numerically true everywhere, with an explicit margin, and unproved.

Two further reductions are recorded for completeness. The rescaling proof of Section~\ref{sec:thm1}
uses relative risk aversion at most one only through a range condition on the true consumption
function of the higher-rate economy, $(1-c_2(0)/c_2(ta))(\rho-1)\le\rho^{1/\mu}-1$ with
$\rho=R_2/R_1$, roughly $c(a)\lesssim\frac{\mu}{\mu-1}c(0)$ (\lean{policy\_mono\_withRate\_crra\_of\_range}).
It is weaker than $\mu\le1$ but fails at moderate wealth. And Theorem 1 holds at equal cash on
hand for all $a$: $g_1(a,z)\le g_2(ta,z)$ with $t=R_1/R_2$ (\lean{policy\_le\_policy\_scaled\_withRate\_of\_marginal\_at}).

\subsection{The invariant-interval form}\label{sec:invariant}

Since Theorem 1 uses the condition only at tomorrow's assets and the coupling only along coupled
paths, if $[0,\bar a']$ is invariant under the lower-rate policy then the condition on $[0,\bar a']$
alone orders the policies there, the coupling started at zero assets never leaves the interval, and
uniqueness follows (\lean{AboveSet}, \lean{push\_aboveSet\_eq\_zero},
\lean{equilibriumRate\_unique\_of\_marginal\_on}). This is the honest form of the reduction: the
condition is needed on the ergodic support and nowhere else. At Aiyagari's numbers the support is
too long for the unweakened condition, as noted; for the slack condition the invariant form is
available but, by Proposition~\ref{prop:homog}, not needed.

\section{Huggett (1993)}\label{sec:huggett}

In Huggett's exchange economy households trade a bond in zero net supply at price $q$; the rate is
$r=1/q-1$, the borrowing limit $\underline a<0$ is a credit limit, and the equilibrium condition is
$K(\mu)=0$. His calibration is $\beta=0.99322$, $\sigma\in\{1.5,3\}$, a two-state endowment, and
credit limits $\underline a\in\{-2,-4,-6,-8\}$ in units of the high endowment. Two things change.

\paragraph{Single crossing needs strictness.}
Against a constant supply, a monotone $K(r)$ gives nothing: two rates with $K=0$ are compatible
with $K\equiv0$ between them. What is needed is that $K$ be strictly increasing, and this follows
if the higher-rate policy dominates the lower-rate one everywhere and strictly on a set the
lower-rate stationary distribution charges (\lean{aggregateCapital\_lt\_of\_policy\_lt},
\lean{huggett\_equilibriumRate\_unique}). The strict form of Theorem 1 holds for creditors: at
$a>0$ with an interior choice, and $r_1<r_2$, $g_1(a,z)<g_2(a,z)$ under \eqref{eq:EC}
(\lean{policy\_lt\_policy\_withRate\_of\_marginal\_at}).

\paragraph{Debtors lose cash on hand when the rate rises.}
For $a<0$, the step $c_2(a)>c_1(a)+(r_2-r_1)a$ has a negative last term, and the chain gives only
\begin{equation}\label{eq:loss}
  g_1(a,z)\ \le\ g_2(a,z)+(r_2-r_1)\max\{0,-a\}:
\end{equation}
saving can fall with the rate for a debtor, but by no more than the extra interest on the debt
(\lean{policy\_le\_policy\_add\_withRate\_of\_marginal\_at}; the slack version is
\lean{policy\_le\_policy\_add\_withRate\_of\_slack\_at}). For $a\ge0$ the loss vanishes and
Theorem 1 is as before.

\paragraph{Carrying the loss through the coupling.}
The coupling of Section~\ref{sec:coupling} needs exact dominance. A version that tolerates a shift
replaces first-order dominance by \emph{dominance up to $\delta$}: $\int h\,d\mu\le\int h\,d\nu+\delta L$
for every test function increasing in assets and $L$-Lipschitz in assets. One period of the
operator turns a shift $\delta$ into $\delta K+\gamma$, where $\gamma$ is the pointwise policy gap
\eqref{eq:loss} and $K$ is a Lipschitz constant of the higher-rate policy in assets; iterating and
passing to the limit,
\begin{equation}\label{eq:shift}
  K(\mu_1)\ \le\ K(\mu_2)+\frac{(r_2-r_1)\,|\underline a|}{1-K}\qquad\text{if }K<1
\end{equation}
(\lean{DominatesShift}, \lean{dominatesShift\_iterate},
\lean{huggett\_aggregateCapital\_le\_add}). Here $K$ is the marginal propensity to save out of
assets. For CRRA the policy's slope in assets is $R(1-\mathrm{MPC})$, and asymptotically the MPC is
$\kappa$, giving $K=R(1-\kappa)=\Thorn=(\beta R)^{1/\sigma}<1$: impatience, with almost no room. At
Huggett's numbers $1-\Thorn_2\approx0.00067$, so the amplification in \eqref{eq:shift} is about
$1500$.

\paragraph{The propensity-to-save bound cannot be had on a bounded region.}
It is tempting to derive the MPC floor from Carroll--Kimball concavity of the consumption function
together with the minimal-MPC level bound $c(x)\ge\kappa x$. On a ray this works: a concave function
above the line $\kappa x$ has every secant slope at least $\kappa$, by one chord far to the right
(\lean{secant\_ge\_of\_concaveOn\_Ici}, with the chord chosen explicitly). On the region
$[\underline a,\bar a]$ it does not: the constant function at the level of the bound at the cap is
concave, lies above the bound, and is flat (\lean{exists\_concaveOn\_ge\_flat}). The MPC floor is an
asymptotic property of the uncapped problem, and the chord at the cap says nothing because there
the level bound is far below actual consumption. So \eqref{eq:shift} is restated with the floor as a
hypothesis in the form it is used: a consumption slope floor $(R_2-K)(b-a)\le c_2(b)-c_2(a)$ is
exactly the Lipschitz bound $K$ for the policy (\lean{policy\_lipschitz\_of\_consumption\_slope}),
and the CRRA form takes $K=\Thorn_2$ and the floor $\kappa_2R_2(b-a)\le c_2(b)-c_2(a)$
(\lean{crra\_huggett\_aggregateCapital\_le\_add\_of\_slack}).

\paragraph{Where Huggett stands.}
Uniqueness reduces to a quantitative race: the creditors' strict gain in \lean{huggett\_equilibriumRate\_unique}
against the debtors' loss bounded by \eqref{eq:shift}. We have the sign of the gain and no lower
bound on its size; every attempt at one runs into a pivot condition of the type \eqref{eq:EC}. And
the MPC floor on the region is the same asymptotic structure that blocks the slack condition at
$\mu>1$. Both are open.

\section{Auxiliary results}\label{sec:aux}

The development contains a number of results that were built as tools for the main theorem or as
tests of the frontier and stand on their own. This section lists the ones an economist may want.

\paragraph{Multiplicity is exactly non-monotone supply, and the equilibrium system permits it.}
Two distinct equilibrium rates force capital supply to be non-monotone on any interval containing
them (\lean{not\_monotoneOn\_of\_equilibria\_ne}), so every multiplicity example in the literature
is a counterexample to Light's Theorem 1 and to nothing else. Conversely an explicit continuous,
strictly positive supply schedule, Cobb--Douglas demand plus a cubic wiggle, has exactly three
equilibria with the demand $D(r)=A^2/(4(r+\delta)^2)$ (\lean{wiggleSupply\_equilibria}),
matching the three-equilibrium shape \citet{YoungWalsh2025} report: nothing in the firm side, the
fixed-point machinery or the intermediate-value argument rules multiplicity out. What is not
proved is that such a schedule is the capital supply of any household; the two bounds available,
the gain test and the minimal MPC, both move the wrong way for a backward bend.

\paragraph{All capital is precautionary.}
A two-state economy near log utility holds at least $1/400$ of capital at every stationary
distribution. The same economy with its income risk removed, both states replaced by their mean,
saves nothing at any asset level, so its capital supply is exactly zero
(\lean{riskless\_aggregateCapital\_eq\_zero}). The corner condition of Section~\ref{sec:corner}
holds at every asset level in the riskless economy and fails above a threshold in the risky one,
and the gap is the value of the insurance the household is buying.

\paragraph{Impatience is necessary.}
\citet{MarcetObiolsWeil2007}, Proposition 3: no stationary equilibrium with finite capital exists at
$\beta R\ge1$. Formalised with income risk through the Euler route, requiring $\beta R>1$ strictly,
and deterministically through secant slopes (\lean{Impatience}); the version with a supply floor is
Theorem~\ref{thm:above}.

\paragraph{Carroll--Kimball concavity across HARA.}
The consumption function is concave in assets for CRRA, shifted CRRA (via a power-mean
inequality), and CARA (via a soft-minimum construction); quadratic utility is proved inadmissible
because it is not increasing (\lean{not\_monotoneOn\_quadUtility}). The induction runs along
Bellman iterates, so it needs the Euler equation one step at a time and at states where the
constraint binds, which is what \lean{IncomeFluctuationEuler} supplies. A witness shows Carroll and
Kimball's conclusion is false when the cap binds (\lean{not\_concaveOn\_consumptionFn\_of\_cap\_binds}),
which is why the cap slack condition cannot be dropped.

\paragraph{A calibrated class with existence and uniqueness.}
Before the log calibration was reached, existence and uniqueness were proved for a class of
economies (\lean{Calibrated}) each of whose numerical fields is one primitive inequality: near-log
CRRA, log Stone--Geary (\lean{stoneGeary\_exists\_equilibrium\_of\_technology'}), and the $b\neq0$
branch of HARA, with general finite Markov income. These carry the machinery of
Section~\ref{sec:arch} at parameters where the published constants do work, and a CES witness
proved uncalibratable shows the cap conditions are not vacuous.

\paragraph{Sufficiency of the Euler equation, without differentiability of the value function.}
A solution of the functional Euler equation with equality at interior choices and the correct
inequality at the constraint is the optimal consumption function
(\lean{eq\_consumptionFn\_of\_isEulerSolution}). The proof is the contraction of the Coleman
operator in the marginal-utility supremum metric of \citet{LiStachurski2014}, run on a class of
candidate policies with a consumption floor and a saving slack. This is the sufficient direction
of the Coleman--Reffett conjugacy of \citet{SargentStachurski2024}, Proposition 8.3.13, proved
without the conjugacy.

\paragraph{Abstract fixed-point results from Sargent and Stachurski.}
Parametric monotonicity of fixed points (Vol.\ 1, Prop.\ 2.2.7); eventually contracting maps
(Thm.\ 6.1.5), which is what spectral-radius conditions deliver; the fitted-value-iteration error
bound (Vol.\ 2, Thm.\ 9.1.3), the value-error input to the policy error bounds of \citet{Li2015};
and uniqueness of fixed points of increasing concave maps, in one dimension and in Du's form on
a finite product (\lean{OrderedFixedPoints}). \citet{Santos2000}'s policy error bound was
assessed and not formalised: it needs interior solutions, which the borrowing constraint denies.

\paragraph{Mean-field-game uniqueness, and why it does not reach Bewley.}
The Lasry--Lions argument for uniqueness of mean-field equilibria adds two optimality inequalities.
Its measure-theoretic core is formalised for the static game and for the finite-horizon
finite-state game on paths: two mean-field Nash distributions have nonpositive Lasry--Lions pairing;
under monotone costs the pairing vanishes; under strict monotonicity the costs coincide, though the
distributions need not (\lean{cost\_eq\_of\_nash}, \lean{value\_eq\_of\_pathEquilibrium}; the
static form of \citealp{CannarsaCapuani2018}, Thm.\ 4.1). The argument cancels
$\int(u_1-u_2)\,d(m_1(0)-m_2(0))$ by a common initial distribution, which a stationary discounted
equilibrium does not have; and \citet{GraberMatter2024} show that games coupled through a market
price are typically not Lasry--Lions monotone, the right condition there being monotonicity of the
error map $Q\mapsto Q-\text{best response}(Q)$, which for Aiyagari is Light's single crossing. The
quantitative version, a strongly monotone Lipschitz operator on a real inner product space has a
unique zero found by a damped iteration with modulus $\sqrt{1-m^2/L^2}$, is formalised
(\lean{exists\_unique\_zero\_of\_stronglyMonotone}); for excess supply it would need a slope floor
and a Lipschitz constant in the rate, neither of which is available.

\paragraph{Tauchen's method.}
The Tauchen transition matrix is defined from cumulative rows of the standard normal distribution
function, taken from Mathlib's Gaussian measure, and proved to be stochastic, positive in its first
column, and monotone in the sense of \ref{E1} for $\rho\ge0$
(\lean{monotoneTransitions\_of\_tauchen}). This is what allows Corollary~\ref{cor:tauchen} to be
stated for any number of grid points.

\section{Conclusion}

Aiyagari's log economy has a unique stationary equilibrium, and the statement is now a theorem
about his model as he wrote it, with his parameters, for any of his earnings processes, checked
by a machine. The proof is the one the picture suggested, supply rises and demand falls, but the
constants in the published versions of that argument do not survive $\beta=0.96$, and what does
survive are consequences of the Euler inequality at the borrowing constraint: a corner condition, a
minimal marginal propensity to consume, and a Doeblin atom. None of it uses a derivative of the
value function.

For $\mu\in\{3,5\}$ we have reduced uniqueness to an inequality between two consumption functions
that holds numerically at every wealth level with an explicit margin and that no available bound
proves; for Huggett, to a quantitative race between creditors and debtors whose creditor side we
cannot bound below. We regard both as genuinely open and have tried to say precisely what a proof
would need: comparative statics of the consumption function in the interest rate, uniform on the
support in relative terms. That is a question about the income fluctuation problem, not about
equilibrium, and it may be the right next target.

On method, the experience of this development is that a proof assistant changes what one writes
down. Arguments that quantify over ``a sufficiently large cap'' or ``a small enough interval of
rates'' had to be given their constants, and at Aiyagari's parameters several of the published
constants were found to be too small by orders of magnitude; the replacements are simpler than the
originals. Differentiability, which enters every textbook treatment through the envelope theorem,
was never needed. And the exact location of the obstruction at $\mu>1$, the wasted interest-income
term in a four-line argument, was found by reading the formal proof rather than the informal one.


\bibliographystyle{ecta}
\bibliography{refs}

\appendix
\section{Theorem index}\label{app:index}

Each result stated in the text, with the Lean declaration that proves it and the file that contains
it. All files are under \lean{LeanEconomics/}. Every declaration listed depends only on the three
standard axioms of Mathlib, \lean{propext}, \lean{Classical.choice} and \lean{Quot.sound}, as
reported by \lean{\#print axioms}; the development contains no \lean{sorry}.

\begin{footnotesize}
\begin{longtable}{>{\raggedright\arraybackslash}p{0.20\textwidth}>{\raggedright\arraybackslash}p{0.43\textwidth}>{\raggedright\arraybackslash}p{0.29\textwidth}}
\toprule
Result & Declaration & File\\
\midrule
\endhead
Theorem~\ref{thm:main} & \lean{aiyagari1994\_log\_equilibriumRate\_unique\_Ioo} & Equilibrium/\allowbreak AiyagariUniqueness\\
Theorem~\ref{thm:main}, interval form & \lean{aiyagari1994\_log\_equilibriumRate\_unique} & Equilibrium/\allowbreak AiyagariUniqueness\\
Theorem~\ref{thm:main}, at the wage & \lean{aiyagari1994\_equilibriumRate\_unique\_wage} & Equilibrium/\allowbreak HomotheticDistribution\\
Corollary~\ref{cor:tauchen} & \lean{aiyagariTauchen\_equilibriumRate\_unique\_Ioo} & Models/\allowbreak AiyagariTauchenWitness\\
Tauchen chain monotone & \lean{monotoneTransitions\_of\_tauchen} & Models/\allowbreak Tauchen\\
Theorem~\ref{thm:above} & \lean{log\_no\_equilibrium\_of\_patient\_risk} & Equilibrium/\allowbreak RateRange\\
Theorem~\ref{thm:above}, CRRA $n$ & \lean{crra\_no\_equilibrium\_of\_patient\_mean} & Equilibrium/\allowbreak Localisation\\
Theorem~\ref{thm:below}(i) & \lean{aiyagari1994\_no\_equilibrium\_of\_budget} & Equilibrium/\allowbreak Localisation\\
Theorem~\ref{thm:below}(ii) & \lean{aiyagari1994\_log\_no\_equilibrium\_below\_mean} & Equilibrium/\allowbreak IncomeMarginal\\
Theorem~\ref{thm:exist} & \lean{aiyagari1994\_log\_exists\_equilibrium\_of\_floor} & Equilibrium/\allowbreak AiyagariExistence\\
Theorem~\ref{thm:bisect} & \lean{aiyagari1994\_log\_bisection} & Equilibrium/\allowbreak Bisection\\
\midrule
Theorem 1, log, local & \lean{crra\_policy\_mono\_withRate\_of\_minMPC} & Models/\allowbreak LargeBetaLight\\
Theorem 1, rescaling core & \lean{monotoneOn\_bellman\_sub\_scaled\_of\_piv} & Models/\allowbreak IncomeFluctuationRateStep\\
Corner from Euler & \lean{minIncome\_le\_consumptionFn\_floor} & Models/\allowbreak EulerCorner\\
Minimal MPC & \lean{crra\_minMPC\_mul\_le\_consumptionFn} & Models/\allowbreak MinimalMPC\\
Cap slack & \lean{crra\_policy\_lt\_cap\_of\_minMPC} & Models/\allowbreak MinimalMPC\\
Log ceiling \eqref{eq:ceiling} & \lean{log\_aggregateCapital\_le} & Equilibrium/\allowbreak AiyagariUniqueness\\
Theorem 2, coupling & \lean{monotoneOn\_capitalSupply\_of\_policy\_mono} & Equilibrium/\allowbreak CapitalSupplyMonotone\\
Doeblin uniqueness & \lean{stationary\_unique} & Distribution/\allowbreak Uniqueness\\
Doeblin convergence & \lean{tendsto\_pushProb\_iterate} & Distribution/\allowbreak Uniqueness\\
Unique stationary distribution, log & \lean{log\_existsUnique\_isStationary} & Equilibrium/\allowbreak AiyagariUniqueness\\
Theorem 3, single crossing & \lean{eq\_of\_monotoneOn\_of\_strictAntiOn} & Equilibrium/\allowbreak Uniqueness\\
Homogeneity of the distribution & \lean{isStationary\_scaleProb\_log} & Equilibrium/\allowbreak HomotheticDistribution\\
Continuity of supply & \lean{continuousOn\_capitalSupply} & Equilibrium/\allowbreak AiyagariContinuity\\
Euler inequalities & \lean{euler\_le}, \lean{euler\_ge} & Models/\allowbreak IncomeFluctuationEuler\\
Income marginal is invariant & \lean{invariant\_eq\_of\_reach} & Equilibrium/\allowbreak IncomeMarginal\\
\midrule
Proposition~\ref{prop:elas} & \lean{policy\_le\_policy\_withRate\_of\_marginal\_at} & Equilibrium/\allowbreak ElasticityReduction\\
CRRA form of \eqref{eq:EC} & \lean{crra\_elasticity\_iff} & Equilibrium/\allowbreak InvariantRegion\\
Uniqueness under \eqref{eq:EC} & \lean{crra\_equilibriumRate\_unique\_of\_marginal} & Equilibrium/\allowbreak ElasticityReduction\\
Invariant-interval form & \lean{crra\_equilibriumRate\_unique\_of\_marginal\_on} & Equilibrium/\allowbreak InvariantRegion\\
Failure at high wealth & \lean{crra\_elasticity\_fails\_of\_asymptotic\_mpc} & Equilibrium/\allowbreak InvariantRegion\\
Proposition~\ref{prop:slack} & \lean{policy\_le\_policy\_withRate\_of\_slack\_at} & Equilibrium/\allowbreak SlackEuler\\
CRRA form of \eqref{eq:SEC} & \lean{crra\_slack\_elasticity\_iff} & Equilibrium/\allowbreak SlackEuler\\
Proposition~\ref{prop:homog} & \lean{crra\_slack\_elasticity\_of\_patience} & Equilibrium/\allowbreak SlackEuler\\
Proposition~\ref{prop:homog}, eventual form & \lean{crra\_slack\_elasticity\_holds\_of\_asymptotic\_mpc} & Equilibrium/\allowbreak SlackEuler\\
Uniqueness under \eqref{eq:SEC} & \lean{crra\_equilibriumRate\_unique\_of\_slack\_on} & Equilibrium/\allowbreak SlackEuler\\
Range condition & \lean{policy\_mono\_withRate\_crra\_of\_range} & Models/\allowbreak IncomeFluctuationRateStep\\
Equal cash on hand & \lean{policy\_le\_policy\_scaled\_withRate\_of\_marginal\_at} & Equilibrium/\allowbreak DebtorRate\\
\midrule
Huggett single crossing & \lean{huggett\_equilibriumRate\_unique} & Equilibrium/\allowbreak Huggett\\
Strict Theorem 1, creditors & \lean{policy\_lt\_policy\_withRate\_of\_marginal\_at} & Equilibrium/\allowbreak DebtorRate\\
Debtor loss \eqref{eq:loss} & \lean{policy\_le\_policy\_add\_withRate\_of\_marginal\_at} & Equilibrium/\allowbreak DebtorRate\\
Shifted dominance \eqref{eq:shift} & \lean{huggett\_aggregateCapital\_le\_add} & Equilibrium/\allowbreak ShiftedDominance\\
Secant slope on a ray & \lean{secant\_ge\_of\_concaveOn\_Ici} & Equilibrium/\allowbreak SavingPropensity\\
Bounded-region obstruction & \lean{exists\_concaveOn\_ge\_flat} & Equilibrium/\allowbreak SavingPropensity\\
MPC floor is a Lipschitz bound & \lean{policy\_lipschitz\_of\_consumption\_slope} & Equilibrium/\allowbreak SavingPropensity\\
Huggett bound, CRRA, slack & \lean{crra\_huggett\_aggregateCapital\_le\_add\_of\_slack} & Equilibrium/\allowbreak SavingPropensity\\
\midrule
Multiplicity forces non-monotone supply & \lean{not\_monotoneOn\_of\_equilibria\_ne} & Equilibrium/\allowbreak Multiplicity\\
Three equilibria & \lean{wiggleSupply\_equilibria} & Equilibrium/\allowbreak Multiplicity\\
Riskless twin saves nothing & \lean{riskless\_aggregateCapital\_eq\_zero} & Models/\allowbreak Precautionary\\
Euler sufficiency & \lean{eq\_consumptionFn\_of\_isEulerSolution} & Models/\allowbreak ColemanReffett\\
Sargent--Stachurski results & \lean{fixedPoint\_le\_of\_dominates}, \lean{tendsto\_iterate\_of\_eventuallyContracting}, \lean{dist\_fixedPoint\_le\_of\_fittedVFI}, \lean{eq\_of\_fixedPt\_of\_concave\_du} & DynamicProgramming/\allowbreak OrderedFixedPoints\\
Lasry--Lions, static & \lean{cost\_eq\_of\_nash} & Equilibrium/\allowbreak LasryLions\\
Lasry--Lions, finite horizon & \lean{value\_eq\_of\_pathEquilibrium} & Equilibrium/\allowbreak LasryLions\\
Strongly monotone operators & \lean{exists\_unique\_zero\_of\_stronglyMonotone} & Analysis/\allowbreak StronglyMonotone\\
Quadratic utility inadmissible & \lean{not\_monotoneOn\_quadUtility} & Models/\allowbreak Quadratic\\
Concavity fails when the cap binds & \lean{not\_concaveOn\_consumptionFn\_of\_cap\_binds} & Models/\allowbreak IncomeFluctuationCarrollKimball\\
\bottomrule
\end{longtable}
\end{footnotesize}

\section{The formalisation}\label{app:formal}

\paragraph{What was checked.}
The development is written in Lean 4 \citep{Lean4}, toolchain \lean{v4.34.0-rc2}, against the
Mathlib library \citep{Mathlib} at a fixed revision recorded in the repository manifest. It
comprises 115 source files, about 35,000 lines, and about 1,500 theorem declarations. A Lean proof
is a term whose type is the statement; the kernel checks that the term is well-formed. What the
reader is asked to trust is the kernel, the axioms, and that the formal statement means what the
prose says. The first two are shared with every Lean development. For the third, each result in the
text names its declaration, and Section~\ref{sec:model} describes the formal objects in enough
detail that the statements can be read; the hypotheses in the text are the hypotheses in the code.

\paragraph{The household problem.}
The dynamic programme is built on a general stochastic dynamic-programming layer with
extended-real rewards, so that $u=-\infty$ at zero consumption is handled without a consumption
floor. The principle of optimality, existence and uniqueness of the bounded value function,
Berge's maximum theorem for the policy (which Mathlib lacks; ours is built on its hemicontinuity
definitions), and the Blackwell contraction are theorems of that layer. The asset region
$[\underline a,\bar a]$ is a type parameter, which is what lets the same theorems serve Aiyagari and
Huggett.

\paragraph{The rate as a state.}
Continuity of the policy in the interest rate is obtained by adjoining the rate to the state and
solving one dynamic programme on the product, so that joint continuity in the state is continuity
in the rate. This avoids any parametric-continuity argument, which would need a modulus that the
Bellman contraction does not obviously supply.

\paragraph{Secant slopes.}
Lipschitz bounds, monotonicity in assets, monotonicity in the rate and the Euler inequalities are
all proved by comparing plans, with no derivative of the value function. Where the development
proves the value function differentiable it uses the sandwich lemma of \citet{ClausenStrub2020},
formalised as \lean{DifferentiableSandwich}. The uniqueness proof does not depend on it.

\paragraph{Probability.}
Stationary distributions are Borel probability measures on a compact metric space, in Mathlib's
\lean{ProbabilityMeasure} type with the weak topology; the Krylov--Bogolyubov existence argument,
weak continuity of aggregation, and the Doeblin contraction on bounded continuous functions are
proved in \lean{Distribution/}. The normal distribution behind Tauchen's method is Mathlib's
\lean{gaussianReal}.

\paragraph{Reproducing the check.}
With Lean and Mathlib installed, \lean{lake build} at the repository root compiles the
development, and \lean{\#print axioms} on any declaration in the index reports the axioms it uses.
The repository is public and the history records the order in which the results were obtained.


\end{document}
